% Author-approved reflection, 22 October 2026; 418 words before TeX markup.
My contribution centred on defining the research requirements and directing the
development and evaluation of the project. I treated the efficient backend as a core
project requirement and kept subsequent development focused on it. I required the
implementation to recover its output from its own maintained state and maintained a
clear distinction between implemented functionality, bounded testing, complexity
arguments, and measured performance. These requirements guided the assessment of
proposed changes and the presentation of results.

I used Codex, Claude, and ChatGPT for technical and language assistance. They
contributed implementation suggestions, code changes, proof formulations, and review
findings. I communicated requirements, set priorities, requested explanations, and
decided which revisions to pursue. When a data structure was unfamiliar, I asked for
examples explaining the relationships between the Jordan algorithm, sibling groups,
and the red--black tree. Particularly useful assistance included finding literature,
discussing the project architecture, expressing ideas academically, and checking LaTeX
syntax. I used ChatGPT's web interface for translation, algorithm explanations, and
discussions of review findings and project planning.

The review process required decisions about what the evidence could support.
Claude-assisted review identified a missing premise in the English proof exposition
and an incomplete diagnostic table. I incorporated these findings into the revision
process while distinguishing corrections to the exposition from changes to the
implementation. The thesis reports that the measured tree configuration had higher
per-input median runtimes than the list backend throughout the formal comparison. I
accepted this outcome and maintained the requirement to report the observed results
and their limitations accurately.

I cannot identify a specific task in this project that I believe would have been
faster without AI. Nevertheless, checking suggestions, understanding unfamiliar
concepts, and revising incomplete answers required effort. The omitted proof premise
illustrates why fluent wording cannot establish that an argument is complete.
Similarly, a successful test run supports a bounded claim about the tested cases; it
does not establish a general correctness or complexity result. These distinctions
remained important when assessing proposed explanations and revisions.

Overall, I am satisfied with the outcome. AI assistance helped me organize the work,
explore technical questions, and express ideas more clearly. The experience reinforced
the importance of clear questions, explicit acceptance criteria, and sufficient
understanding to assess proposed changes. In future work, I would document my
decisions and the reasons for accepting or rejecting suggestions more systematically,
alongside the AI contributions themselves. This would make both my reasoning and the
role of the tools easier to trace. Responsibility for understanding and defending the
submitted work remains with me.
